\documentclass[aps,prx,reprint,floatfix,superscriptaddress]{revtex4-2}
\usepackage[T1]{fontenc}
\usepackage[utf8]{inputenc}
\usepackage{microtype}
\usepackage{amssymb,amsthm,mathtools}
\usepackage{bm}
\usepackage{booktabs}
\usepackage{array}
\usepackage{graphicx}
\usepackage{xcolor}
\usepackage{hyperref}
\usepackage[nameinlink,noabbrev]{cleveref}
\hypersetup{hidelinks}

% revtex4-2 4.2f and array 2.6n compatibility (TeX Live 2026).
\makeatletter
\let\switch@array\relax
\makeatother

\newtheorem{definition}{Definition}
\newtheorem{theorem}{Theorem}
\newtheorem{corollary}{Corollary}
\newtheorem{proposition}{Proposition}
\newtheorem{lemma}{Lemma}
\newtheorem{remark}{Remark}

\newcommand{\DataNumber}[2]{#2}
\newcommand{\Ftwo}{\mathbb F_2}
% Generated by scripts/build_all_tables.py; do not edit.
\newcommand{\TradeoffRuns}{\DataNumber{w3-cells}{592}}
\newcommand{\MatchedCells}{\DataNumber{w4-cells}{1,080}}
\newcommand{\MatchedCompleted}{\DataNumber{w4-completed}{1,065}}
\newcommand{\MatchedPredicateErrors}{\DataNumber{w4-predicate-errors}{15}}
\newcommand{\ExternalCells}{\DataNumber{w5-cells}{1,728}}
\newcommand{\ExternalCompleted}{\DataNumber{w5-completed}{1,695}}
\newcommand{\ExternalPredicateErrors}{\DataNumber{w5-predicate-errors}{33}}
\newcommand{\CertificateMutationCases}{\DataNumber{w6-mutations}{56}}
\newcommand{\CertificateFalseAccepts}{\DataNumber{w6-false-accepts}{0}}
\newcommand{\CertificateCrosscheckCases}{\DataNumber{w6-crosschecks}{12}}
\newcommand{\QRECells}{\DataNumber{w7-cells}{288}}
\newcommand{\QREMaterializeSpacetimeRatio}{\DataNumber{w7-spacetime-ratio}{2814}}
\newcommand{\NaturalCells}{\DataNumber{w8-cells}{21,060}}
\newcommand{\NaturalStableFamilies}{\DataNumber{w8-stable-families}{5}}
\newcommand{\NaturalNegativeRegressions}{\DataNumber{w8-negative-regressions}{0}}
\newcommand{\PresetGateEffect}{\DataNumber{w8-preset-effect}{364.9}}
\newcommand{\PresetGateEffectLow}{\DataNumber{w8-preset-effect-low}{349.8}}
\newcommand{\PresetGateEffectHigh}{\DataNumber{w8-preset-effect-high}{380.5}}
\newcommand{\UnifiedRows}{\DataNumber{w9-unified-rows}{24,748}}
\newcommand{\UnifiedPredicateErrors}{\DataNumber{w9-predicate-errors}{48}}
\newcommand{\CleanTheoryTests}{\DataNumber{w9-clean-theorem-tests}{9}}
\newcommand{\CleanGridPairs}{\DataNumber{w9-clean-grid-pairs}{17}}


\begin{document}

\title{Representation-Constrained Recoverability in Quantum Compilation:\\
A $K$-ary Randomized Multipass Frontier and Certified Measurements}

\author{Jinze Yang}
\affiliation{School of Physics, Xidian University, Xi'an 710071, China}

\author{Yangyang Li}
\email{yyli@xidian.edu.cn}
\affiliation{Key Laboratory of Intelligent Perception and Image Understanding of Ministry of Education of China,
Joint International Research Laboratory of Intelligent Perception and Computation,
International Research Center for Intelligent Perception and Computation,
Collaborative Innovation Center of Quantum Information of Shaanxi Province,
School of Artificial Intelligence, Xidian University, Xi'an 710071, China}

\author{Xiu-Hao Deng}
\email{dengxiuhao@iqasz.cn}
\affiliation{Shenzhen International Quantum Academy, Shenzhen 518048, China}
\affiliation{Shenzhen Branch, Hefei National Laboratory, Shenzhen 518048, China}

\date{August 3, 2026}

\begin{abstract}
Equivalent quantum circuits can expose different amounts of recoverable
structure to a compiler.  We make this representation dependence explicit
with a cyclic $K$-ary family and a dispersed, round-major additive-update
stream in which no token reveals the final coefficient.  Exact wraparound
analysis gives an $\epsilon$-packing of size
$\exp[\Theta(m\log(1/\epsilon))]$.  For randomized compilers making
$p=p(m)$ forward passes, failing with probability at most $\delta$, committing
$A_p$ prefix-output bits, and carrying complete serialized crossing state of
worst-input expected size $S$, we prove
\(
 \overline A_p+(2p-1)S\geq
 (1-\delta)m\log_2K-h_2(\delta)
\), together with separate family-relative and self-contained output bounds.
A memory-capped block compiler reaches the transcript envelope within a
constant factor, while a semantic aggregate reaches the compact-output
corner.  We then measure the formal resources directly, compare matched
representations and semantic-capable external baselines, certify every
completed output, and propagate a common total error budget through an
end-to-end fault-tolerant estimator.  The natural-workload campaign contains
\NaturalCells\ certified cells; \NaturalStableFamilies\ structured families
satisfy a predeclared all-size/all-seed criterion, while the two negative
controls show \NaturalNegativeRegressions\ systematic regressions.  The
result is a recoverability frontier and a reproducible validation package,
not a claim that one compiler representation dominates every circuit.
\end{abstract}

\maketitle

\section{Question, scope, and claims}
\label{sec:introduction}

Compiler pipelines routinely move between circuit syntax, Pauli networks,
phase polynomials, and ZX-like global intermediate representations.  Such
choices matter because aggregation that is explicit in one representation
may have to be inferred from a long stream in another.  Phase-polynomial and
Pauli-gadget methods exploit this structure before returning to gates
\cite{amy2014tdepth,nam2018automated,cowtan2020phase}; ZX methods provide a
global language in which phases can propagate and merge
\cite{duncan2020graph,kissinger2020reducing,kissinger2020pyzx}.  Full-stack
resource studies then amplify upstream choices into physical cost
\cite{fellous2023optimizing,watkins2024surfacecode}.

We ask a narrower, operational question: how much information must a
compiler either commit early or carry across an input cut in order to recover
one member of an approximately separated unitary family?  The answer depends
on the input representation and on what the resource model charges.  Our
model therefore serializes all control, store, window, parameter, and live-IR
state.  RAM, caches, random seeds, dynamic pass state, and rereadable
materialized input cannot be hidden outside the cut budget.

The paper has four claims, each paired with evidence.  First, the packing and
update representations implement the same target unitary and have exact
projective separation.  Second, the randomized multipass lower bound follows
from conditional decoding and Fano's inequality.  Third, explicit semantic
and block-hybrid compilers achieve the stated upper corners, with no stronger
self-contained gate-count optimality claim.  Fourth, instrumented,
certificate-gated experiments occupy the permitted theorem region and expose
when the constructive advantage persists or disappears.  Legacy witnesses,
stress tests, and complete provenance tables are separated into the
Supplemental Material.

\section{Recoverability frontier}
\label{sec:theory}

\subsection{Distance, family, and charged compiler model}

For equal-dimension unitaries define the projective operator distance
\begin{equation}
 d_{\rm proj}(U,V):=\min_{|\omega|=1}\lVert U-\omega V\rVert_{\rm op}.
 \label{eq:projective-distance}
\end{equation}
The output code is fixed before the input is drawn and is prefix-free.  A
\emph{family-relative} code may use a public dictionary of generators and
instance dimensions; a \emph{self-contained} code must serialize all
addresses, parameters, and circuit framing.  We write their output lengths
as $B_{\rm rel}^{\rm out}$ and $B_{\rm self}^{\rm out}$, respectively.

\begin{definition}[Independent-character packing base]
\label{def:packing-construction}
Let $P_1,\ldots,P_m$ be commuting Pauli-$Z$ characters whose joint sign map
onto $\{\pm1\}^m$ is surjective; $P_j=Z_j$ on $m$ qubits is the canonical
choice.  The public family dictionary fixes this ordered list and the
address codec.
\end{definition}

At a named input cut, the restart state is the self-delimiting concatenation
\begin{equation}
 B_{\rm cut}=B_{\rm control}+B_{\rm store}+B_{\rm window}
             +B_{\rm parameter}+B_{\rm IR}.
 \label{eq:info-budget}
\end{equation}
The fields include all mutable compiler state needed to resume exactly.
Output bytes become charged commitments when they are irrevocably appended.
The family-relative hybrid output $E_{\rm rel}$ uses a public mode tag and
block boundaries; unaggregated coordinates remain addressed update records.

\subsection{Exact cyclic $K$-ary packing}
\label{sec:w1-qary-packing}

The no-wrap construction of \cref{def:packing-construction} is convenient
for output counting, but modular additive shares require a grid whose period
is also the sharing modulus.  The following formulation records exactly when
wraparound creates a collision.

\begin{definition}[Wrap margin]
\label{def:w1-wrap-margin}
For an alphabet $[K]_0$, step $\Delta>0$, and $K\geq2$, define
\begin{align}
 \operatorname{dist}_{2\pi}(a,0)
 &:=\min_{z\in\mathbb Z}|a-2\pi z|,\\
 \gamma(K,\Delta)
 &:=\min_{1\leq d\leq K-1}
       \operatorname{dist}_{2\pi}(d\Delta,0).
\end{align}
Let $P_1,\ldots,P_m$ be independent Pauli-$Z$ characters as in
\cref{def:packing-construction}, and set
\begin{equation}
 U_x^{(K,\Delta)}:=
 \exp\!\left[-\frac{i}{2}\sum_{j=1}^m x_j\Delta P_j\right],
 \qquad x\in[K]_0^m.
\end{equation}
\end{definition}

\begin{theorem}[Exact worst-pair separation]
\label{thm:w1-exact-separation}
For every $m\geq1$,
\begin{equation}
 \min_{x\neq y}d_{\mathrm{proj}}
 \bigl(U_x^{(K,\Delta)},U_y^{(K,\Delta)}\bigr)
 =2\sin\!\left(\frac{\gamma(K,\Delta)}4\right).
 \label{eq:w1-exact-separation}
\end{equation}
Consequently the family is projectively injective exactly when
$\gamma(K,\Delta)>0$.  If $(K-1)\Delta\leq\pi$, then
$\gamma(K,\Delta)=\Delta$.  If $Q\geq K+1$ and
$\Delta=2\pi/Q$, then again $\gamma(K,\Delta)=\Delta$, even though the
coefficient interval $[0,(K-1)\Delta]$ may cross the principal-argument
cut at $\pi$.
\end{theorem}

\begin{proof}
Fix $x\neq y$, put $s=x-y$, and choose $j$ with $s_j\neq0$.
Independence of the characters makes their joint sign map surjective.
Hence the relative unitary $U_x^{(K,\Delta)}
(U_y^{(K,\Delta)})^\dagger$ has two eigenvalues obtained from sign patterns
that agree off $j$ and whose circular angular separation is
$\operatorname{dist}_{2\pi}(|s_j|\Delta,0)\geq\gamma(K,\Delta)$.
After any common phase rotation, one of two circle points separated by
$\eta\in[0,\pi]$ is at chordal distance at least
$2\sin(\eta/4)$ from $1$.  This proves the lower bound in
\eqref{eq:w1-exact-separation} after taking the global-phase infimum.

Choose $d\in\{1,\ldots,K-1\}$ attaining $\gamma(K,\Delta)$ and take a pair that
differs by $d$ in one coordinate only.  Its relative spectrum consists of
two points with circular separation $\gamma(K,\Delta)$, so centering the
shorter arc attains $2\sin(\gamma/4)$.  This proves equality.  Under the
no-wrap hypothesis, every $d\Delta$ lies in $[\Delta,\pi]$, whose minimum
is $\Delta$.  For $\Delta=2\pi/Q$ and
$1\leq d\leq K-1\leq Q-2$, every nonzero grid distance is at least
$2\pi/Q$, with equality at $d=1$.
\end{proof}

\begin{corollary}[$\epsilon$-tuned cyclic family]
\label{cor:w1-epsilon-cyclic}
For $0<\epsilon\leq1/8$, let
\begin{equation}
 Q_\epsilon:=\left\lfloor
 \frac{\pi}{2\arcsin(2\epsilon)}\right\rfloor,
 \qquad K_\epsilon:=Q_\epsilon-1,
 \qquad \Delta_\epsilon:=\frac{2\pi}{Q_\epsilon}.
 \label{eq:w1-epsilon-grid}
\end{equation}
Then
\begin{equation}
 \frac1{4\epsilon}\leq K_\epsilon\leq\frac{\pi}{4\epsilon},
 \quad
 \min_{x\neq y}d_{\mathrm{proj}}(U_x,U_y)
 =2\sin\!\left(\frac{\pi}{2Q_\epsilon}\right)
 \geq4\epsilon.
 \label{eq:w1-epsilon-pack}
\end{equation}
Thus, with $\mathcal G_{m,\epsilon}:=
\{U_x:x\in[K_\epsilon]_0^m\}$,
\begin{equation}
 \log_2|\mathcal G_{m,\epsilon}|
 =m\log_2K_\epsilon
 =\Theta\!\left(m\log\frac1\epsilon\right).
 \label{eq:w1-packing-entropy}
\end{equation}
The allowed coefficient interval is
$[0,(Q_\epsilon-2)2\pi/Q_\epsilon]$.  It can extend beyond $\pi$ but
omits the residue $Q_\epsilon-1$; \cref{eq:w1-exact-separation}, rather
than a principal-interval argument, controls its wraparound.
\end{corollary}

\begin{proof}
The bound $\arcsin(2\epsilon)\leq\pi\epsilon$ and
$\lfloor a\rfloor\geq a-1$ give
$K_\epsilon\geq1/(2\epsilon)-2\geq1/(4\epsilon)$; the reverse bound uses
$\arcsin(2\epsilon)\geq2\epsilon$.  Since
$Q_\epsilon\leq\pi/[2\arcsin(2\epsilon)]$,
$\pi/(2Q_\epsilon)\geq\arcsin(2\epsilon)$.  Apply
\cref{thm:w1-exact-separation}.  The cardinality statement is immediate.
\end{proof}

\subsection{Dispersed additive-update representation}
\label{sec:w1-dispersed}

\begin{definition}[$r$-round dispersed update stream]
\label{def:w1-dispersed-stream}
The cyclic specialization of \cref{def:phase-family} with modulus
$Q\geq3$, $K=Q-1$, $\Delta=2\pi/Q$, $r\geq2$ rounds, and
$\Sigma=\{+1\}$: for $x\in[K]_0^m$ the shares
$a^{(1)},\ldots,a^{(r)}\in\mathbb Z_Q^m$ satisfy
\begin{equation}
 \sum_{t=1}^r a_j^{(t)}=x_j\pmod Q
 \qquad(j\in[m]),
 \label{eq:w1-share-sum}
\end{equation}
the flat input is the round-major record sequence
$\mathsf{Update}(t,j,a_j^{(t)})$, and the semantic input is the addressed
aggregate vector $x$.
\end{definition}

\begin{proposition}[Semantic equivalence and token privacy]
\label{prop:w1-stream-equivalence}
Every valid flat stream of \cref{def:w1-dispersed-stream} implements the
same $U_x$ as the semantic input, up to global phase.  Moreover, a single
flat token never determines $x_j$: for every fixed $(t,j,a)$ and every
$b\in[K]_0$, there is a valid stream with $a_j^{(t)}=a$ and $x_j=b$.
Under the canonical sharing distribution in which the first $r-1$ shares
are independent uniform residues and the last is fixed by
\eqref{eq:w1-share-sum}, every individual residue is uniform on
$\mathbb Z_Q$ and independent of $x_j$.
\end{proposition}

\begin{proof}
For integer representatives there are $k_j\in\mathbb Z$ with
$\sum_ta_j^{(t)}=x_j+Qk_j$.  Commutativity and $Q\Delta=2\pi$ give
\begin{equation*}
 \prod_tR_{P_j}(\Delta a_j^{(t)})
 =R_{P_j}(\Delta x_j+2\pi k_j)
 =(-1)^{k_j}R_{P_j}(\Delta x_j).
\end{equation*}
The product of the signs is global.  For privacy, fix the named token and
choose any one of the other $r-1$ residues to enforce the desired sum;
the remaining residues are arbitrary.  The distributional claim follows
because a sum containing at least one independent uniform residue is
uniform, including the last share after conditioning on $x_j$.
\end{proof}

\begin{definition}[Forward-pass cut accounting]
\label{def:w1-pass-accounting}
Partition the round-major stream into Alice's prefix $A$ (rounds
$1,\ldots,r-1$) and Bob's suffix $B$ (round $r$).  A forward $p$-pass
compiler scans $AB$ from left to right at most $p=p(m)$ times.  There are
$p$ crossings $A\to B$ and at most $p-1$ rewind crossings $B\to A$.
At every crossing, the complete restart state is the five-field
serialization of \eqref{eq:info-budget}; rereadable input and every live
IR object are therefore included.  For a run on $I$, let $A_p(I)$ be the
number of bits
irreversibly committed to the final output while the machine scans $A$,
summed over all passes, and define
\begin{equation}
 \Sigma_p(I):=\sum_{t=1}^{p}B_{\mathrm{cut}}(c_{A\to B,t})
       +\sum_{t=1}^{p-1}B_{\mathrm{cut}}(c_{B\to A,t}).
 \label{eq:w1-transcript-budget}
\end{equation}
Both are literal, self-delimiting transcript lengths.  Unbarred resource
symbols take maxima over valid inputs.  For randomized compilers, write
\begin{equation*}
 \overline A_p:=\max_I\mathbb E_\rho A_p(I,\rho),\qquad
 \overline\Sigma_p:=\max_I\mathbb E_\rho\Sigma_p(I,\rho),
\end{equation*}
and let
$S:=\max_I\mathbb E_\rho[\max_cB_{\mathrm{cut}}(c;I,\rho)]$ over the
$2p-1$ named crossings.  Then
$\overline\Sigma_p\leq(2p-1)S$; the deterministic version is identical
without expectations.
\end{definition}

\begin{theorem}[One-pass zero-error mask bound]
\label{thm:w1-one-pass-mask}
Let a deterministic one-pass compiler be correct to
$\epsilon<\sin(\pi/(2Q))$ on every valid stream of
\cref{def:w1-dispersed-stream}.  Then
\begin{equation}
 A_1+B_{\mathrm{cut}}(c_{A\to B,1})
 \geq\left\lceil m\log_2Q\right\rceil.
 \label{eq:w1-one-pass-mask}
\end{equation}
\end{theorem}

\begin{proof}
For each $u\in\mathbb Z_Q^m$, choose a canonical prefix whose aggregate
is $u$ (put $u$ in its first round and zero in the other prefix rounds).
The committed prefix and restart state after this prefix must distinguish
all $Q^m$ choices.  Indeed, for $u\neq u'$, choose one common final share
$v$ coordinatewise.  At each coordinate at most two residues of $v_j$
would make either $u_j+v_j$ or $u'_j+v_j$ equal to the single excluded
aggregate $Q-1$; since $Q\geq3$, an allowed $v_j$ exists.  Both completions
are valid, and they differ in every coordinate in which $u$ and $u'$ differ.
By \cref{thm:w1-exact-separation}, one circuit cannot be within
$\epsilon$ of both targets.  The messages are prefix-free under the
declared serializers, so Kraft's inequality gives
\eqref{eq:w1-one-pass-mask}.
\end{proof}

\subsection{Randomized multipass frontier}
\label{sec:w1-randomized-multipass}

Put
\begin{equation}
 L_\delta(m,K):=(1-\delta)m\log_2K-h_2(\delta),
 \qquad 0\leq\delta<\tfrac12.
 \label{eq:w1-Ldelta}
\end{equation}

\begin{theorem}[$K$-ary randomized multipass recoverability frontier]
\label{thm:w1-unified-frontier}
Let $Q\geq3$, $K=Q-1$, and
$0\leq\epsilon<\sin(\pi/(2Q))$.  Let a possibly randomized compiler be
correct within projective error $\epsilon$, with probability at least
$1-\delta$ for every $x\in[K]_0^m$ and every valid $r$-round dispersed
representation of $x$.  Assume that it makes at most $p=p(m)$ forward
passes in the charged model of \cref{def:w1-pass-accounting}.  Then:
\begin{align}
 \overline B_{\mathrm{rel}}^{\mathrm{out}}
 &\geq L_\delta(m,K),
 &
 \overline B_{\mathrm{self}}^{\mathrm{out}}
 &\geq L_\delta(m,K),
 \label{eq:w1-output-lower}\\
 \overline A_p+\overline\Sigma_p
 &\geq L_\delta(m,K).
 \label{eq:w1-pass-lower}
\end{align}
In particular, for the worst-input expected crossing cap $S$ of
\cref{def:w1-pass-accounting},
\begin{equation}
 \boxed{\quad
 \overline A_p+(2p-1)S
 \geq(1-\delta)m\log_2K-h_2(\delta).
 \quad}
 \label{eq:w1-frontier-box}
\end{equation}
Here $S$ includes RAM, parameters, cache/hash metadata, materialized input,
and all IR fields from \eqref{eq:info-budget}; it is not a RAM-only bound.
For the tuned choice of \cref{cor:w1-epsilon-cyclic}, both right-hand sides
are $\Omega((1-\delta)m\log(1/\epsilon))$ for fixed
$\delta<1/2$.
\end{theorem}

\begin{proof}
For \eqref{eq:w1-output-lower}, take $X$ uniform on $[K]_0^m$ and fix one
valid dispersed representation for each value of $X$.
The packing separation is greater than $2\epsilon$, so nearest-packing
decoding recovers $X$ from a successful output.  Fano's inequality gives
\begin{equation*}
 I(X;C)\geq m\log_2K-h_2(\delta)
 -\delta\log_2(K^m-1)\geq L_\delta(m,K).
\end{equation*}
Expected length in either separately fixed prefix-free output code is at
least this mutual information; a worst-input expectation is at least the
uniform average.

For \eqref{eq:w1-pass-lower}, independently choose a uniform prefix
aggregate $U\in\mathbb Z_Q^m$, put it in the first prefix round, put zero
in any other prefix rounds, and set the last share $V=X-U\pmod Q$.
Then $V$ is uniform and independent of $X$, so
$H(X\mid V)=m\log_2K$.  Convert the $p$ scans into the standard alternating
simulation across the $A\mid B$ partition.  The transcript $T$ contains
each complete crossing snapshot and all output chunks written on Alice's
side.  Given $T$ and $V$, Bob can reproduce every Bob-side segment and the
complete final output; random-tape state and dynamic control are present in
the crossing serialization.  The output decoder therefore recovers $X$
with error at most $\delta$, and conditional Fano yields
\begin{equation*}
 I(X;T\mid V)\geq L_\delta(m,K).
\end{equation*}
The framed transcript is prefix-free conditional on $V$, so its expected
literal length is at least this information.  Its length is precisely
$A_p+\Sigma_p$ up to fields already included in the self-delimiting
serializers.  Maximizing over inputs proves \eqref{eq:w1-pass-lower}; the
cap and tuned-$\epsilon$ consequences follow.
\end{proof}

\begin{remark}[What is and is not optimal]
\label{rem:w1-optimality-scope}
Equation~\eqref{eq:w1-output-lower} is an output-information bound and
\eqref{eq:w1-frontier-box} is a transcript bound; neither converts
self-contained circuit bits into gates.  The factor $2p-1$ counts the
actual alternating crossings of a forward scan.  A random-access input
oracle, an uncharged rereadable file, or reverse scans would be a different
model.  For $p=1$ and zero error, \cref{thm:w1-one-pass-mask} strengthens
$m\log_2K$ to $m\log_2(K+1)$.  No such strengthening is claimed for
randomized or general multipass protocols.
\end{remark}

\subsection{Memory-capped constructions}
\label{sec:w1-upper}

\begin{proposition}[Block hybrid]
\label{prop:w1-block-hybrid}
Fix $p,h\geq1$ and let $q=\min\{m,ph\}$.  There is a deterministic
zero-error $p$-pass compiler for \cref{def:w1-dispersed-stream} that
aggregates $q$ coordinates and passes the other update records through in
the public $p$-block hybrid mode of $E_{\mathrm{rel}}$.  It obeys
\begin{align}
 S&\leq \left\lceil\frac qp\right\rceil\log_2Q
       +O(\log(mQrp)),
 \label{eq:w1-upper-state}\\
 A_p&\leq(r-1)(m-q)\log_2Q+O(r\log(mQrp)),
 \label{eq:w1-upper-A}\\
 B_{\mathrm{rel}}^{\mathrm{out}}
 &\leq q\log_2K+r(m-q)\log_2Q
       +O(r\log(mQrp)).
 \label{eq:w1-upper-output}
\end{align}
In the compact-output corner $h=\lceil m/p\rceil$, this becomes
\begin{align}
 B_{\mathrm{rel}}^{\mathrm{out}}
 &\leq m\log_2K+O(\log(mQrp)),\\
 S&\leq\left\lceil\frac mp\right\rceil\log_2Q
       +O(\log(mQrp)).
 \label{eq:w1-compact-upper}
\end{align}
For two rounds, the construction satisfies
$A_p+pS\leq m\log_2Q+O(p\log(mQrp))$ after using a smaller final block.
Compared with
\eqref{eq:w1-frontier-box} at $\delta=0$, it is within the constant
$(2p-1)\log_2(K+1)/(p\log_2K)<2\log_2 3$ for every $K\geq2$.
\end{proposition}

\begin{proof}
Partition the first $q$ coordinates into at most $p$ consecutive blocks of
size at most $h$.  On pass $t$, keep only the residues of block $t$ and add
their first $r-1$ rounds modulo $Q$; when the final-round mates arrive,
emit the corresponding aggregate base-$K$ block.  The live packed block,
indices, and framing give \eqref{eq:w1-upper-state}.  On the first pass,
emit every share of each unselected coordinate in the declared round-major
order; these rotations already have the correct product and require no
retained coefficient table.  The selected coordinates can be balanced
across the $p$ blocks, so the maximum block has size $\lceil q/p\rceil$.
Their consecutive boundaries are a public function of $(m,p,q)$, so the
header stores those integers rather than $p$ separate delimiters.
The prefix contains $r-1$ residues and the
whole output contains $r$ residues for each such coordinate, giving
\eqref{eq:w1-upper-A} and \eqref{eq:w1-upper-output}.  If $ph\geq m$ there
are no passthrough coordinates.  For $r=2$, $(m-q)\log_2Q+ph\log_2Q$ is
at most $m\log_2Q$ up to the last partial block and framing.  The constant
comparison uses $(2p-1)/p<2$ and
$\log_2(K+1)/\log_2K\leq\log_2 3$.
\end{proof}

\begin{proposition}[Semantic and randomized upper corners]
\label{prop:w1-semantic-upper}
Under the same public dictionary, a one-pass semantic compiler receiving
$x$ emits the packed base-$K$ aggregate code with
\begin{equation}
 B_{\mathrm{rel}}^{\mathrm{out}}
 \leq\lceil m\log_2K\rceil+O(\log(mK)),
 \qquad S=O(\log(mK)),
 \label{eq:w1-semantic-upper}
\end{equation}
and implements the same $U_x$ as the flat stream.  For any
$0\leq\delta<1/2$, an input-independent abort coin that runs this exact
compiler with probability $1-\delta$ and otherwise emits a fixed short
circuit has worst-input success at least $1-\delta$ and expected relative
output
\begin{equation}
 (1-\delta)\lceil m\log_2K\rceil+O(\log(mK)+1).
 \label{eq:w1-random-upper}
\end{equation}
Thus the randomized family-relative output lower bound is matched up to
framing and the binary-entropy additive term.  A self-contained
$R_Z$-circuit upper bound remains
$O(m(\log m+\log Q))$ bits and at most $m$ rotations; no matching
self-contained-bit or gate-count claim is made.
\end{proposition}

\begin{proof}
Pack the ordered aggregate vector as one base-$K$ integer and retain only
its streaming index and framing state.  Proposition~\ref{prop:w1-stream-equivalence}
gives semantic equality with the flat representation.  Prefix the output
with the abort-coin branch tag for the randomized construction and take
expectations.  The self-contained construction writes each nonzero
$R_{P_j}(2\pi x_j/Q)$ with its address and rational parameter.
\end{proof}


\subsection{Interpretation of the bound}

The lower bound is an information statement, not a gate-count statement.
It holds for every valid dispersed representation, including the canonical
sharing distribution, and permits $p$ to grow with $m$.  The $2p-1$ factor is
the number of named forward and rewind crossings in the declared scan model.
A reverse-scanning compiler, random-access oracle, or uncharged input file is
a different model.  At the compact family-relative corner the block hybrid
uses $S=O((m/p)\log(K+1))$; the semantic compiler receives $x$ directly and
needs only logarithmic framing state.  This difference is the precise
representation dependence studied here.

\section{Constructions and certified measurements}
\label{sec:experiments}

\subsection{Certificate gate}

Every completed result entering a quality mean is independently certified.
The exact symbolic domain and soundness proposition are stated below; dense
projective matrices provide a second checker through six qubits.  Unsupported,
predicate-error, and numerical-inconclusive outcomes remain status records and
are never converted into scores.

% W6: exact symbolic correctness certificate and soundness domain.
% This file is input by PaperDraft/ucc_v7_2_prxq_repaired/main.tex.

Let


\begin{equation*}
 \mathcal C_n=\langle H,S,S^\dagger,X,Y,Z,\mathrm{CX},
 \mathrm{CZ},\mathrm{SWAP}\rangle
\end{equation*}
be the Clifford gate language on addressed qubits.  Fix a declared finite
formal-symbol set $\mathcal S$ and the exact coefficient module
\begin{equation*}
 \Lambda=\mathbb Q\pi\oplus
 \bigoplus_{s\in\mathcal S}\mathbb Q\theta_s.
\end{equation*}
For a signed Hermitian $n$-qubit Pauli $P$ and $\lambda\in\Lambda$, write
$R_P(\lambda)=\exp(-i\lambda P/2)$.

\begin{definition}[Exact commuting-Pauli/Clifford certificate domain]
\label{def:certificate}
The domain $\mathcal D_{\rm CPC}$ consists of finite addressed circuits over
$\mathcal C_n$, exact records $R_Z(\lambda)$ and $R_P(\lambda)$, identity and
barrier records, and an arbitrary exact global phase.  Every coefficient is
a reduced rational $\pi$ component plus a sorted sparse list of reduced
rational formal-symbol components carried by the input; a floating-only
parameter is not in the exact domain.

Scan a circuit in execution order while maintaining a Clifford frame $F$ and
the invariant
\begin{equation}
 U_{\rm prefix}=F R_{P_t}(\lambda_t)\cdots R_{P_1}(\lambda_1)
 \quad\text{up to global phase}.
 \label{eq:w6-frame-invariant}
\end{equation}
The frame is represented by the complete signed inverse-conjugation tableau
$\{F^\dagger X_jF,F^\dagger Z_jF\}_{j=1}^n$.  Appending a physical rotation
$R_Q(\lambda)$ appends $F^\dagger QF$ to the Pauli list.  The certificate
predicate additionally requires all extracted $P_j$ to commute pairwise.

Canonicalization moves a Pauli sign into its coefficient, discards an
identity-axis rotation, adds coefficients on identical unsigned supports,
reduces the rational $\pi$ coefficient modulo $2$ to the interval $[0,2)$,
keeps every formal-symbol coefficient exact, deletes zero records, and sorts the remaining records lexicographically by
their binary $(x,z)$ masks.  Circuit and Clifford global phases are omitted.
The resulting record is
\begin{align}
 \beta_j&:=(b_{j,s})_{s\in\mathcal S}\in\mathbb Q^{\mathcal S},
 &\bar a_j&\in\mathbb Q\cap[0,2),\nonumber\\
 \operatorname{can}(C)&=
 \bigl(n,T_F,((x_j,z_j,\bar a_j,\beta_j))_{j=1}^{\ell}\bigr).
 \label{eq:w6-canonical-record}
\end{align}
The symbolic checker accepts exactly when the two canonical byte records are
equal.
\end{definition}

\begin{proposition}[Soundness of canonical symbolic acceptance]
\label{prop:symbolic-soundness}
For $C,C'\in\mathcal D_{\rm CPC}$ satisfying the pairwise-commutation
predicate,
\begin{equation*}
 \operatorname{can}(C)=\operatorname{can}(C')
 \quad\Longrightarrow\quad
 U_C=e^{i\phi}U_{C'}
\end{equation*}
for some $\phi\in\mathbb R$.
\end{proposition}

\begin{proof}
Induction over the input records proves \eqref{eq:w6-frame-invariant}.
For a Clifford record $G$, replace $F$ by $GF$.  For a rotation,
\begin{equation*}
 R_Q(\lambda)F=F R_{F^\dagger QF}(\lambda),
\end{equation*}
which is exactly the tableau update used by the checker.  Thus the extracted
record implements the original circuit up to its explicitly ignored global
phase.

Pairwise commutation permits deterministic reordering without changing the
unitary.  Sign normalization is exact because
$R_{-P}(\lambda)=R_P(-\lambda)$, while coefficient reduction changes a factor only
by
\begin{equation*}
 R_P(\lambda+2k\pi)=(-1)^kR_P(\lambda).
\end{equation*}
All such signs multiply to one input-independent global phase; an
identity-axis rotation is also a global phase.  Equality of the sorted
rotation records, including every formal-symbol component, therefore gives
equal rotation products projectively.

Finally, equality of the signed tableaus means that
$F'F^\dagger$ commutes with every Pauli generator.  The Pauli representation
generates the full matrix algebra, so $F'F^\dagger$ is a scalar.  Hence the
two Clifford frames also agree up to global phase, proving the claim.
\end{proof}

\begin{remark}[Status and completeness boundary]
\label{rem:w6-certificate-status}
Only status \texttt{completed\_valid}, accompanied by the content-addressed
canonical records and an independent certificate payload, is an acceptance.
A canonical mismatch is \texttt{completed\_invalid}; an unlisted gate or
malformed address is \texttt{unsupported}; failure of pairwise commutation is
\texttt{predicate\_error}; and a floating coefficient without an exact
sidecar, an unassigned formal symbol in the dense checker, or a dense distance
inside the numerical guard band is
\texttt{numerical\_inconclusive}.  Only \texttt{completed\_valid} rows enter
quality means.  The proposition is a soundness statement, not a claim that
every equivalent Clifford-plus-rotation circuit has the same record.
\end{remark}

For $n\leq6$, the audit independently forms dense matrices and checks the
normalized Frobenius projective distance
\begin{equation*}
 d_{\rm F,proj}(U,V)=2^{-n/2}\min_{|\omega|=1}\|U-\omega V\|_{\rm F}.
\end{equation*}
The frozen suite covers both accepted and rejected pairs at every width
$1\leq n\leq6$ and changes, separately, an angle, a sign, one support bit,
a qubit permutation, the Clifford frame, and a deleted or duplicated gate.
The executable definition and status contract are
\path{certificates/CERTIFICATE_SPEC.md}; the canonicalizer and checker are
\path{certificates/canonicalize.py} and
\path{certificates/symbolic_checker.py}.


The frozen audit contains \CertificateMutationCases\ adversarial mutations
and \CertificateCrosscheckCases\ dense/symbolic cross-checks.  It observed
\CertificateFalseAccepts\ false accepts and no checker disagreements.  This
is finite executable evidence for the stated domain, not a completeness proof
outside it.

\subsection{Direct tradeoff tracing}

The reference suite implements echo, full aggregation, three memory-capped
hybrids, and an IR-materializing control.  Each event records committed output
bytes and the five cut fields in Eq.~\eqref{eq:info-budget}, separately from
RSS.  Across \TradeoffRuns\ frozen runs, every measured point satisfies the
theorem diagnostic; no point was deleted as an outlier.

\begin{figure*}[t]
 \centering
 \includegraphics[width=0.93\textwidth]{../../figures/measured_tradeoff_frontier.pdf}
 \caption{Measured commitment--crossing-state--IR surface from event traces.
 The theoretical lower envelope is evaluated in the same serialized-bit
 coordinates.  RSS is recorded but is not substituted for $S$.}
 \label{fig:tradeoff}
\end{figure*}

\subsection{Matched representations and strong external baselines}

The factorial benchmark holds target, target basis, hardware model, timeout,
memory cap, total error budget, and certificate fixed while changing only the
representation and compiler.  It contains \MatchedCells\ cells across nine
representations and five reference compilers; \MatchedCompleted\ completed
and the remaining fifteen are retained predicate errors.  The external
campaign contains \ExternalCells\ cells covering local Qiskit, semantic-capable
TKET PauliSimp and PyZX full-reduce, and the authors' minimal phase-table
reference, each in native and unified-basis configurations.  Of these,
\ExternalCompleted\ completed; bridge loss is reported separately.

\begin{figure*}[t]
 \centering
 \includegraphics[width=0.95\textwidth]{../../figures/representation_compiler_interaction.pdf}
 \caption{Measured representation--compiler interaction.  Completion status
 is shown separately from metrics; predicate errors are not numerical wins.}
 \label{fig:matched}
\end{figure*}

The external results change the interpretation: a tool that recovers a phase
polynomial before synthesis belongs on the semantic-capable side, regardless
of authorship.  The evidence therefore supports a representation-sensitive
comparison, not an authors-versus-tools ranking.

\subsection{Fixed-total-error fault-tolerant QRE}

All pipelines use the same additive split
$\epsilon_{\rm total}=\epsilon_{\rm alg}+\epsilon_{\rm synth}
+\epsilon_{\rm logical}$ and allocate the synthesis share across the actual
rotation multiset.  The \QRECells\-cell campaign crosses
$\epsilon_{\rm total}\in\{10^{-3},10^{-6},10^{-9}\}$, two surface-code
profiles, multiple factory-parallelism settings, equal and optimized
allocation policies, and semantic-first, materialize-first, and external
semantic-capable pipelines.  It reports rotations, logical $T$ states, data
and factory qubits, factories, cycles, runtime, and spacetime volume.

\begin{table*}[t]
\caption{Representative fixed-error slice: diagonal witness,
$\epsilon_{\rm total}=10^{-6}$, conservative profile, equal allocation, and
sixteen requested factories.  This is one registered slice, not the entire
sensitivity campaign.}
\label{tab:qre}
\centering
\begin{tabular}{lrrrrrr}
\toprule
Pipeline & Rotations & $T$ states & Physical qubits & Factories & Cycles & Qubit-seconds\\
\midrule
% Generated from data/frozen/qre_summary.csv.
External semantic-capable (TKET) & 29 & 798 & 258,152 & 16 & 13,317 & 3.44e+03 \\
Materialize-first & 23,992 & 2,679,110 & 468,968 & 16 & 57,966,528 & 2.72e+07 \\
Semantic-first reference & 22 & 1,800 & 305,000 & 16 & 31,675 & 9.66e+03 \\
\bottomrule

\end{tabular}
\end{table*}

On this slice, materialize-first uses \QREMaterializeSpacetimeRatio$\times$
the spacetime of the authors' semantic-first reference.  TKET also recovers a
compact semantic form and is therefore reported on the semantic-capable side.
The ordering is stable over the frozen grid where the ratio remains above
one; exact physical values remain conditional on the two declared QEC models.

\begin{figure*}[t]
 \centering
 \includegraphics[width=0.47\textwidth]{../../figures/qre_pareto.pdf}\hfill
 \includegraphics[width=0.47\textwidth]{../../figures/qre_sensitivity.pdf}
 \caption{End-to-end QRE Pareto surface (left) and sensitivity across total
 error, QEC profile, allocation, and factory parallelism (right).}
 \label{fig:qre}
\end{figure*}

\subsection{Natural workloads, negative controls, and ablation}

The W8 design crosses seven structured families---QAOA/Ising, diagonal
Hamiltonians, commuting-block Trotterization, QPE controlled powers, QFT
arithmetic, controlled-power ladders, and phase-polynomial blocks---plus
random Clifford+$R_Z$ and routing-dominated negative controls.  Three size
levels, three seeds, two densities, and all-to-all/line topologies are crossed
with a complete six-toggle authors-method factorial and the strongest available
external semantic baseline.

\begin{figure*}[t]
 \centering
 \includegraphics[width=0.95\textwidth]{../../figures/natural_workload_summary.pdf}
 \caption{Matched natural-family comparison and negative controls.  Error bars
 are paired bootstrap 95\% intervals; all families and controls are retained.}
 \label{fig:natural}
\end{figure*}

\begin{table}[b]
\caption{External-minus-authors gate difference and predeclared strict
all-cell success.  Positive differences favor the authors reference.}
\label{tab:natural}
\centering
\small
\begin{tabular}{lcc}
\toprule
Family & Mean [95\% interval] & Success\\
\midrule
% Generated from data/frozen/natural_workload_summary.csv.
controlled\_power\_ladder & 9.0 [5.7, 12.7] & yes \\
diagonal\_hamiltonian & 38.0 [25.9, 51.8] & yes \\
phase\_polynomial\_blocks & -24.3 [-33.8, -16.3] & no \\
qaoa\_ising & 44.7 [33.7, 56.6] & yes \\
qft\_arithmetic & 8.1 [1.9, 13.8] & no \\
qpe\_controlled\_powers & -12.0 [-22.0, -3.0] & yes \\
trotter\_commuting\_blocks & 33.0 [17.6, 48.2] & yes \\
\bottomrule

\end{tabular}
\end{table}

Five named families meet the strict criterion, but QFT arithmetic and the
phase-polynomial family do not; both remain visible in
Fig.~\ref{fig:natural} and Table~\ref{tab:natural}.  The two negative controls
show \NaturalNegativeRegressions\ systematic quality regressions.  The largest
isolated gate-count contribution is the preset recognizer: a mean reduction of
\PresetGateEffect\ gates with bootstrap interval
[\PresetGateEffectLow, \PresetGateEffectHigh].  Selector, cache/reuse, and
projected-block selection have zero isolated gate-count effect in this
reference implementation, although their runtime effects and all pairwise
interactions are retained.

\begin{figure*}[t]
 \centering
 \includegraphics[width=0.93\textwidth]{../../figures/ablation_effects.pdf}
 \caption{All six factorial main effects and pairwise gate interactions.
 Zero-effect factors are shown rather than omitted.}
 \label{fig:ablation}
\end{figure*}

\section{Interpretation, limitations, and conclusion}
\label{sec:discussion}

The theorem identifies a recoverability constraint: in the dispersed flat
representation, aggregate information must appear in committed output or in
the charged crossing transcript.  It does not say that a gate list is always
hard, that semantic structure cannot be recovered, or that an authors'
compiler must beat a global external optimizer.  Indeed, TKET and PyZX are
explicitly classified as semantic-capable whenever they recover global phase
structure before synthesis.

The experiments are controlled reference measurements.  Dense certification
limits the natural factorial to six qubits; increasing scale repeats and
densifies structure rather than increasing Hilbert-space width.  Runtime and
RSS are machine-specific.  The QRE is end-to-end under two versioned
surface-code assumptions, but it is not an architecture-independent forecast.
The block-hybrid upper bound is constant-factor matching for the declared
family-relative transcript model; no matching self-contained bit bound or
universal gate-count optimum is claimed.

Data handling is part of the result.  A normalized index contains
\UnifiedRows\ rows with unique nine-field primary keys.  It retains
\UnifiedPredicateErrors\ predicate-error rows outside quality means.  Every
experimental number printed in this manuscript is generated from a registered
frozen source, and the one-command rebuild runs data audits, theorem tests,
benchmark and QRE subsets, full tests, both TeX builds, and reference checks.

In summary, representation is a quantifiable compilation resource when the
input, output code, passes, and complete state channels are declared.  The
$K$-ary randomized multipass theorem gives the lower frontier; the semantic
and memory-capped compilers give attainable corners; and the certified
measurements show both stable gains and deliberate negative results.  This
combination supports representation-aware benchmarking while keeping claims
inside the proved model.

\section*{Data and code availability}

The canonical data registry is \path{data/manifest.yaml}, the normalized
schema is \path{schema/experiment.schema.json}, and the numeric provenance map
is \path{data/provenance_map.yaml}.  Full experiment reports, raw logs,
certificates, QRE configurations, hashes, and clean-room receipts accompany
the repository.  Supplemental Material contains extended statuses, legacy
witnesses, stress results, and the claim--evidence matrix.

\bibliography{references}

\end{document}
